% Copyright (c) 2022 by Lars Spreng
% This work is licensed under the Creative Commons Attribution 4.0 International License.

\documentclass[
11pt,notheorems,hyperref={pdfauthor=whatever}
]{beamer}

\usepackage{kotex}
\input{loadslides.tex}

\usepackage{fontspec}
\IfFileExists{/Users/seojunha/Library/Fonts/NotoSansCJKkr-Regular.otf}{%
  \setmainhangulfont[
    Path=/Users/seojunha/Library/Fonts/,
    BoldFont=NotoSansCJKkr-Bold.otf
  ]{NotoSansCJKkr-Regular.otf}
  \setsanshangulfont[
    Path=/Users/seojunha/Library/Fonts/,
    BoldFont=NotoSansCJKkr-Bold.otf
  ]{NotoSansCJKkr-Regular.otf}
  \setmonohangulfont[
    Path=/Users/seojunha/Library/Fonts/,
    BoldFont=NotoSansCJKkr-Bold.otf
  ]{NotoSansCJKkr-Regular.otf}
}{%
  \IfFontExistsTF{Noto Sans CJK KR}{%
    \setmainhangulfont{Noto Sans CJK KR}
    \setsanshangulfont{Noto Sans CJK KR}
    \setmonohangulfont{Noto Sans CJK KR}
  }{%
    \IfFontExistsTF{NanumGothic}{%
      \setmainhangulfont{NanumGothic}
      \setsanshangulfont{NanumGothic}
      \setmonohangulfont{NanumGothic}
    }{}%
  }%
}
\renewcommand{\familydefault}{\sfdefault}

\usepackage{tikz}
\usetikzlibrary{arrows.meta, positioning, calc, fit, backgrounds, shapes.geometric}
\usepackage{booktabs}
\usepackage{array}

\definecolor{TextGray}{HTML}{333333}
\definecolor{DeepBlue}{HTML}{08457E}
\definecolor{SoftBlue}{HTML}{EAF0F8}
\definecolor{LineGray}{HTML}{B8B8B8}
\definecolor{WarnRed}{HTML}{B3261E}
\definecolor{GoodGreen}{HTML}{237A4B}
\definecolor{SoftGreen}{HTML}{EAF6EF}
\definecolor{SoftRed}{HTML}{FBEAEA}
\definecolor{SoftYellow}{HTML}{FFF4D8}
\definecolor{SoftPurple}{HTML}{F1EAFE}

\addtobeamertemplate{frametitle}{}{\vspace{-0.5em}}
\setbeamerfont{title}{size=\LARGE,series=\bfseries}
\setbeamerfont{subtitle}{size=\normalsize}
\setbeamerfont{frametitle}{series=\bfseries}

\tikzset{
  state/.style={draw=DeepBlue, very thick, rounded corners=2pt, fill=SoftBlue, align=center, minimum width=1.55cm, minimum height=0.82cm, inner sep=4pt, font=\scriptsize\bfseries},
  action/.style={draw=LineGray, rounded corners=2pt, fill=white, align=center, minimum width=1.2cm, minimum height=0.50cm, inner sep=3pt, font=\scriptsize},
  value/.style={draw=GoodGreen, rounded corners=2pt, fill=SoftGreen, align=center, minimum width=1.4cm, minimum height=0.50cm, inner sep=3pt, font=\scriptsize\bfseries},
  lossbox/.style={draw=WarnRed, very thick, rounded corners=2pt, fill=SoftRed, align=center, minimum width=2.0cm, minimum height=0.55cm, inner sep=3pt, font=\scriptsize\bfseries},
  bluebox/.style={draw=DeepBlue, very thick, rounded corners=2pt, fill=SoftBlue, align=center, minimum height=0.82cm, inner sep=5pt, font=\small\bfseries},
  greenbox/.style={draw=GoodGreen, very thick, rounded corners=2pt, fill=SoftGreen, align=center, minimum height=0.82cm, inner sep=5pt, font=\small\bfseries},
  redbox/.style={draw=WarnRed, very thick, rounded corners=2pt, fill=SoftRed, align=center, minimum height=0.82cm, inner sep=5pt, font=\small\bfseries},
  yellowbox/.style={draw=orange!70!black, very thick, rounded corners=2pt, fill=SoftYellow, align=center, minimum height=0.82cm, inner sep=5pt, font=\small\bfseries},
  purplebox/.style={draw=purple!70!black, very thick, rounded corners=2pt, fill=SoftPurple, align=center, minimum height=0.82cm, inner sep=5pt, font=\small\bfseries},
  arrow/.style={-{Latex[length=2.5mm]}, thick, draw=DeepBlue},
  grayarrow/.style={-{Latex[length=2.2mm]}, thick, draw=LineGray},
  redarrow/.style={-{Latex[length=2.2mm]}, thick, draw=WarnRed},
  greenarrow/.style={-{Latex[length=2.2mm]}, thick, draw=GoodGreen}
}

\newcommand{\smallgray}[1]{{\scriptsize\textcolor{gray!75}{#1}}}
\newcommand{\term}[1]{\textbf{\textcolor{DeepBlue}{#1}}}

\title[]{연구 아이디어}
\subtitle{}
\author[]{하서준}
\institute{DGIST}
\date{2026년 6월 12일}

\begin{document}

{
\setbeamertemplate{footline}{}
\begin{frame}
  \titlepage
  \vspace{-0.8em}
  \begin{center}
  \scriptsize
  LA-CDM \(\cdot\) LDTL \(\cdot\) 시간적 학습 신호 불일치 \(\cdot\) 미래 인식 단계별 학습
  \end{center}
\end{frame}
}
\addtocounter{framenumber}{-1}

\begin{frame}{기존 방법 1: LA-CDM}
\vspace{0.1em}
\begin{block}{LA-CDM의 학습 방식}
\small
LA-CDM은 HA가 진단 가설과 확신도를 예측하고, DA가 검사 또는 최종 진단을 선택하는 두 에이전트 반복 구조를 학습한다. DA는 진단 경로 단위의 최종 보상을 이용해 GRPO로 학습된다.
\end{block}

\vspace{0.15em}
\begin{center}
\begin{tikzpicture}[scale=0.87, transform shape]
\node[state] (s0) at (0,0) {\(s_0\)};
\node[state] (s1) at (1.75,0) {\(s_1\)};
\node[state] (s2) at (3.50,0) {\(s_2\)};
\node[state] (s3) at (5.25,0) {\(s_3\)};
\node[state] (sf) at (7.00,0) {최종};
\node[action] at (0.88,0.55) {\(a_0\)};
\node[action] at (2.63,0.55) {\(a_1\)};
\node[action] at (4.38,0.55) {\(a_2\)};
\node[action] at (6.13,0.55) {진단};
\draw[arrow] (s0) -- (s1);
\draw[arrow] (s1) -- (s2);
\draw[arrow] (s2) -- (s3);
\draw[arrow] (s3) -- (sf);
\draw[draw=WarnRed, very thick, rounded corners=7pt]
  (0.15,-0.48) -- (0.15,-0.82) -- (6.95,-0.82) -- (6.95,-0.48);
\node[lossbox, minimum width=6.8cm] (r) at (3.55,-1.28)
  {하나의 최종 보상 \(A_{\mathrm{traj}}\)이 전체 행동에 동일 적용};
\end{tikzpicture}
\end{center}

\vspace{0.1em}
\begin{alertblock}{한계}
\small
DA GRPO는 하나의 최종 보상을 경로 내부의 모든 DA 행동 토큰에 동일하게 적용한다.
\begin{itemize}
  \item 최종적으로 맞춘 trajectory에서는 \textbf{진단 개선에 크게 기여한 검사}와 \textbf{우연히 뒤따라온 덜 유용한 검사}가 같은 방향으로 강화될 수 있다.
  \item 반대로 최종적으로 틀린 trajectory에서는 \textbf{초기에 유용한 검사를 선택했더라도} 함께 패널티를 받을 수 있다.
\end{itemize}
\end{alertblock}
\end{frame}

\begin{frame}{기존 방법 2: LDTL}
\vspace{0.1em}
\begin{block}{LDTL의 학습 방식}
\small
LDTL은 단계별 의사결정 학습 신호를 제공하려는 방향에 가깝다. 즉, 행동을 최종 진단 경로 보상 하나로만 보지 않고, 현재 단계의 진단 변화와 연결하려고 한다.
\end{block}

\vspace{0.2em}
\begin{center}
\begin{tikzpicture}[scale=0.95, transform shape]
\node[state] (st) at (0,0) {\(s_t\)};
\node[action] (at) at (2.0,0.55) {\(a_t\)};
\node[state] (st1) at (4.0,0) {\(s_{t+1}\)};
\node[state] (st2) at (6.4,0) {\(s_{t+2}\)};
\node[state] (st3) at (8.8,0) {\(s_{t+3}\)};
\draw[arrow] (st) -- (st1);
\draw[grayarrow] (st1) -- (st2);
\draw[grayarrow] (st2) -- (st3);
\draw[redarrow, bend left=30] (st.north) to node[above, font=\scriptsize] {한 단계 신호} (st1.north);
\node[redbox, minimum width=4.0cm] at (6.6,-1.0) {효과가 뒤늦게 나타날 수 있음};
\draw[redarrow] (6.6,-0.72) -- (st3.south);
\end{tikzpicture}
\end{center}

\vspace{0.2em}
\begin{alertblock}{문제점}
\small
한 단계 뒤 변화만 보면, 특정 검사가 \textbf{바로 확신도를 올리지는 않지만 이후 검사/진단을 가능하게 만드는 지연된 진단 가치}를 충분히 반영하기 어렵다.
\end{alertblock}
\end{frame}

\begin{frame}{공통 문제: 시간적 학습 신호 불일치}
\vspace{0.2em}
\begin{columns}[T,onlytextwidth]
\begin{column}{0.48\textwidth}
\begin{center}
\begin{tikzpicture}
\node[redbox, minimum width=0.94\linewidth] {
LA-CDM\\
\small 행동별 기여도 구분이 없음
};
\end{tikzpicture}
\end{center}
\end{column}
\begin{column}{0.48\textwidth}
\begin{center}
\begin{tikzpicture}
\node[yellowbox, minimum width=0.94\linewidth] {
LDTL\\
\small 행동별 신호는 있지만 한 단계 뒤 변화에 의존함
};
\end{tikzpicture}
\end{center}
\end{column}
\end{columns}

\begin{center}
\vspace{0.85em}
{\large\bfseries\textcolor{DeepBlue}{공통적인 문제}}\\[0.65em]
\small
현재 선택이 앞으로 어떤 영향을 미칠지 모델링하지 못한다.\\
즉, 각 검사 선택이 이후 진단 정확도와 confidence를 얼마나 개선하는지
행동 단위로 학습하기 어렵다.
\end{center}
\end{frame}

\begin{frame}{생각 중인 두 가지 방향}
\vspace{0.25em}
\begin{columns}[T,onlytextwidth]
\begin{column}{0.48\textwidth}
\begin{center}
\begin{tikzpicture}
\node[bluebox, minimum width=0.95\linewidth] {
방향 A\\
\normalsize LA-CDM에 단계별 가중치 부여\\
\small 행동별 미래 인식 보상 할당
};
\end{tikzpicture}
\end{center}
\vspace{0.2em}
\scriptsize
\begin{itemize}
  \item 기존 LA-CDM 구조는 유지.
  \item 최종 보상을 모든 행동에 동일 적용하지 않음.
  \item \(n\)-단계 뒤 진단 개선량으로 행동별 보상 계산.
  \item 단계별 가중치를 통해 행동별 손실을 다르게 적용.
\end{itemize}
\end{column}
\begin{column}{0.48\textwidth}
\begin{center}
\begin{tikzpicture}
\node[greenbox, minimum width=0.95\linewidth] {
방향 B\\
\normalsize 사전학습 진단 에이전트 + Planner\\
\small DA 평가 기반 2-step Planner 학습
};
\end{tikzpicture}
\end{center}
\vspace{0.2em}
\scriptsize
\begin{itemize}
  \item DA는 진단을 잘 맞추고 confidence를 잘 내도록 학습.
  \item 후보 action을 DA에 2단계까지 넣어 미래 가치 평가.
  \item 미래 가치를 가장 높이는 첫 action으로 Planner 학습.
  \item 추론 시 Planner가 현재 state에서 다음 action 선택.
\end{itemize}
\end{column}
\end{columns}

\end{frame}

\begin{frame}{방향 A: LA-CDM에 단계별 가중치 부여}
\vspace{0.15em}
\begin{block}{핵심 아이디어}
\small
LA-CDM의 생성 경로는 그대로 유지하되, 모든 행동에 같은 최종 보상을 주지 않는다. 각 행동이 이후 진단 상태를 얼마나 개선했는지에 따라 \textbf{행동별로 다른 학습 신호}를 준다.
\end{block}

\vspace{0.15em}
\begin{center}
\begin{tikzpicture}[scale=0.88, transform shape]
\node[state] (s0) at (0,0) {\(s_0\)\\\(h_0,c_0\)};
\node[state] (s1) at (1.9,0) {\(s_1\)\\\(h_1,c_1\)};
\node[state] (s2) at (3.8,0) {\(s_2\)\\\(h_2,c_2\)};
\node[state] (s3) at (5.7,0) {\(s_3\)\\\(h_3,c_3\)};
\node[state] (s4) at (7.6,0) {\(s_4\)\\최종};
\node[action] at (0.95,0.55) {\(a_0\)};
\node[action] at (2.85,0.55) {\(a_1\)};
\node[action] at (4.75,0.55) {\(a_2\)};
\node[action] at (6.65,0.55) {\(a_3\)};
\draw[arrow] (s0) -- (s1);
\draw[arrow] (s1) -- (s2);
\draw[arrow] (s2) -- (s3);
\draw[arrow] (s3) -- (s4);
\draw[redarrow, bend left=35] (s0.north) to node[above, font=\scriptsize] {미래 가치 평가} (s2.north);
\draw[redarrow, bend right=35] (s1.south) to node[below, font=\scriptsize] {미래 가치 평가} (s3.south);
\draw[greenarrow, bend left=35] (s2.north) to node[above, font=\scriptsize] {미래 가치 평가} (s4.north);
\end{tikzpicture}
\end{center}

\vspace{0.15em}
\begin{block}{직관적인 절차}
\small
\begin{enumerate}
  \item 한 환자에 대해 기존 LA-CDM처럼 여러 진단 경로를 생성한다.
  \item 각 행동이 이후 상태의 진단 가치를 얼마나 올렸는지 확인한다.
  \item 도움이 된 행동의 토큰 구간에는 더 큰 가중치, 덜 도움이 된 행동에는 더 작은 가중치를 준다.
\end{enumerate}
\end{block}

\vspace{0.05em}
\begin{alertblock}{핵심 메시지}
\small
최종 보상은 유지하되, 그 보상을 진단 경로 전체에 똑같이 뿌리지 않고 행동별 기여도에 따라 다시 나눈다.
\end{alertblock}
\end{frame}

\begin{frame}{방향 A: 왜 단계별 가중치가 필요한가?}
\vspace{0.35em}
\begin{center}
\begin{tikzpicture}[scale=0.86, transform shape]
\draw[gray!45, line width=0.7pt] (-0.2,-1.15) -- (10.0,-1.15);

\node[anchor=east, font=\bfseries\small, text=WarnRed] at (-0.35,0.55) {Before};
\node[state] (bs0) at (0.55,0.55) {\(s_0\)};
\node[state] (bs1) at (2.20,0.55) {\(s_1\)};
\node[state] (bs2) at (3.85,0.55) {\(s_2\)};
\node[state] (bs3) at (5.50,0.55) {\(s_3\)};
\node[state] (bs4) at (7.15,0.55) {최종};
\node[action] at (1.38,1.05) {\(a_0\)};
\node[action] at (3.03,1.05) {\(a_1\)};
\node[action] at (4.68,1.05) {\(a_2\)};
\node[action] at (6.33,1.05) {진단};
\draw[arrow] (bs0) -- (bs1);
\draw[arrow] (bs1) -- (bs2);
\draw[arrow] (bs2) -- (bs3);
\draw[arrow] (bs3) -- (bs4);
\draw[draw=WarnRed, very thick, rounded corners=6pt]
  (0.65,-0.05) -- (0.65,-0.38) -- (7.05,-0.38) -- (7.05,-0.05);
\node[lossbox, minimum width=6.0cm] at (3.85,-0.78)
  {최종 보상 하나가 모든 행동에 동일 적용};

\node[anchor=east, font=\bfseries\small, text=GoodGreen] at (-0.35,-2.9) {After};
\node[state] (as0) at (0.55,-2.9) {\(s_0\)};
\node[state] (as1) at (2.20,-2.9) {\(s_1\)};
\node[state] (as2) at (3.85,-2.9) {\(s_2\)};
\node[state] (as3) at (5.50,-2.9) {\(s_3\)};
\node[state] (as4) at (7.15,-2.9) {최종};
\node[action] (aa0) at (1.38,-2.40) {\(a_0\)};
\node[action] (aa1) at (3.03,-2.40) {\(a_1\)};
\node[action] (aa2) at (4.68,-2.40) {\(a_2\)};
\node[action] (aa3) at (6.33,-2.40) {진단};
\draw[arrow] (as0) -- (as1);
\draw[arrow] (as1) -- (as2);
\draw[arrow] (as2) -- (as3);
\draw[arrow] (as3) -- (as4);
\draw[greenarrow, bend left=35] (as0.north) to node[above, font=\scriptsize] {미래 가치 개선} (as2.north);
\draw[greenarrow, bend right=35] (as1.south) to node[below, font=\scriptsize] {미래 가치 개선} (as3.south);
\draw[greenarrow, bend left=35] (as2.north) to node[above, font=\scriptsize] {미래 가치 개선} (as4.north);
\node[greenbox, minimum width=1.45cm] at (1.38,-3.95) {\(w_0\)};
\node[greenbox, minimum width=1.45cm] at (3.03,-3.95) {\(w_1\)};
\node[greenbox, minimum width=1.45cm] at (4.68,-3.95) {\(w_2\)};
\node[greenbox, minimum width=1.45cm] at (6.33,-3.95) {\(w_3\)};
\node[anchor=west, align=left, font=\scriptsize] at (7.65,-3.95)
  {행동별 기여도에 따라\\서로 다른 학습 신호};
\end{tikzpicture}
\end{center}

\vspace{0.25em}
\begin{block}{핵심}
\small
LA-CDM은 trajectory 내부에서 \textbf{어떤 action이 좋았고 나빴는지}를 모델링하지 못한다.
우리는 action 단위 가중치를 통해 \textbf{각 행동의 기여도를 supervise}한다.
\end{block}
\end{frame}

\begin{frame}{방향 B: Lookahead 학습}
\vspace{0.15em}
\begin{center}
\begin{tikzpicture}[scale=1.02, transform shape]
\draw[gray!45, line width=0.8pt] (4.7,-1.8) -- (4.7,2.0);

\node[font=\bfseries\small, text=DeepBlue] at (2.2,1.75) {Stage 1};
\node[font=\small, text=TextGray] at (2.2,1.45) {Confidence Calibration};
\node[align=right, font=\small, text=TextGray, text width=3.0cm, anchor=east] at (-0.25,0.35)
  {진단을 잘 하고\\confidence를 정확히\\내도록 학습};
\node[greenbox, minimum width=3.8cm, minimum height=1.55cm] (diag) at (2.2,0.35)
  {Diagnosis Agent};

\node[font=\bfseries\small, text=DeepBlue] at (7.2,1.75) {Stage 2};
\node[font=\small, text=TextGray] at (7.2,1.45) {Planner Supervised Training};
\node[bluebox, minimum width=3.8cm, minimum height=1.55cm] (plan) at (7.2,0.35)
  {Planner Agent};
\node[align=left, font=\small, text=TextGray, text width=3.7cm, anchor=west] at (9.30,0.20)
  {Diagnosis Agent의\\confidence 기반\\2-step 평가로\\액션 선택 학습};
\end{tikzpicture}
\end{center}

\vspace{0.25em}
\begin{block}{핵심 아이디어}
\small
먼저 DA는 진단 정확도와 confidence를 학습한다. 이후 Planner는 후보 action을 2-step으로 넣어 보며, 미래 가치를 가장 많이 올리는 action을 학습한다.
\end{block}

\end{frame}

\begin{frame}{방향 B: 2-step Lookahead}
\vspace{0.15em}
\begin{center}
\begin{tikzpicture}[scale=0.86, transform shape]
\node[state, font=\large\bfseries] (s) at (0,0) {\(s_t\)};

\node[greenbox, minimum width=1.15cm] (ct) at (2.0,1.15) {1단계\\CT};
\node[action] (cbc) at (2.0,-1.15) {1단계\\CBC};

\node[action] (ctcbc) at (4.3,1.85) {2단계\\CBC};
\node[action] (ctus) at (4.3,0.65) {2단계\\US};
\node[action] (cbcct) at (4.3,-0.65) {2단계\\CT};
\node[action] (cbcus) at (4.3,-1.85) {2단계\\US};

\node[value] (v1) at (6.5,1.85) {\(V=0.82\)};
\node[value] (v2) at (6.5,0.65) {\(V=0.64\)};
\node[value] (v3) at (6.5,-0.65) {\(V=0.51\)};
\node[value] (v4) at (6.5,-1.85) {\(V=0.57\)};

\node[greenbox, minimum width=2.2cm] (best) at (8.65,1.85) {CT 선택};

\draw[arrow] (s) -- (ct);
\draw[arrow] (s) -- (cbc);
\draw[arrow] (ct) -- (ctcbc);
\draw[grayarrow] (ct) -- (ctus);
\draw[grayarrow] (cbc) -- (cbcct);
\draw[grayarrow] (cbc) -- (cbcus);
\draw[arrow] (ctcbc) -- (v1);
\draw[grayarrow] (ctus) -- (v2);
\draw[grayarrow] (cbcct) -- (v3);
\draw[grayarrow] (cbcus) -- (v4);
\draw[greenarrow] (v1) -- (best);

\node[redbox, minimum width=5.0cm] at (4.2,-3.0)
{DA가 평가한 미래 가치로 첫 행동 선택};
\end{tikzpicture}
\end{center}

\vspace{0.25em}
\begin{block}{Lookahead 학습 신호}
\small
학습 시에는 후보 action을 2단계까지 가정하고, 각 미래 상태를 학습된 DA가 평가한다. Planner는 미래 가치를 가장 많이 올리는 첫 action을 학습한다.
\end{block}

\end{frame}

\begin{frame}{두 방향 비교}
\vspace{0.1em}
\begin{center}
\scriptsize
\begin{tabular}{p{0.22\textwidth}p{0.34\textwidth}p{0.34\textwidth}}
\toprule
 & \textbf{방향 A} & \textbf{방향 B} \\
\midrule
핵심 아이디어 &
LA-CDM에 행동별 \(n\)-단계 미래 인식 가중치 추가 &
DA를 진단 정확도와 confidence 중심으로 학습한 뒤, DA 평가값으로 Planner를 지도학습 \\
\addlinespace
공통 문제 해결 방식 &
최종 보상을 행동별 기여도에 따라 재가중 &
후보 action의 2-step 미래 가치를 비교해 다음 행동 선택 \\
\addlinespace
구현 비용 &
낮음: 기존 LA-CDM의 보상/손실 수정 &
높음: Lookahead 교사 행동 생성과 Planner 학습 필요 \\
\bottomrule
\end{tabular}
\end{center}

\vspace{0.35em}
\begin{block}{공통점}
\centering
\large
현재 action에 미래 가치를 반영해서,\\
이후 진단 상태를 더 잘 개선하는 action 선택을 학습하도록 유도한다.
\end{block}
\end{frame}

\end{document}
